\section{Korai világháló}
A világháló indulása a CERN-hez kötődik: Tim Berners-Lee 1989 márciusában írta meg az első webes javaslatot, majd 1990 végére megvalósította a HTML, a HTTP és az URL alapjait, valamint az első webszervert és böngésző/szerkesztő programot \cite{cern2024www35,w3cLittleWebHistory}. Az első weboldal és információs szerver az info.cern.ch címen futott, egy CERN-ben használt NeXT számítógépen \cite{w3cLittleWebHistory}. Ezek a korai oldalak főként statikus dokumentumok voltak: a webszerver a tárolt HTML-fájlokat küldte vissza a böngészőnek, dinamikus tartalomgenerálás nélkül \cite{mdnWebServer}. A grafikus böngészők közül az NCSA Mosaic 1993-as megjelenése, majd a Mosaic Communications Corp. (később Netscape) 1994-es megalakulása jelentős szerepet játszott a web népszerűsödésében \cite{w3cLittleWebHistory,nsf2004mosaic}. Ezzel párhuzamosan a HTML is folyamatosan fejlődött: a HTML 4.0 1997-ben, a HTML5 pedig 2014-ben lett W3C Recommendation, miközben a HTML alapvető szerepe továbbra is a webes dokumentumok szerkezetének leírása maradt \cite{w3c1997html40,w3c2014html5}.
